\documentclass[11pt]{article}

\usepackage[margin=1in]{geometry}
\usepackage{amsmath,amssymb,amsthm}
\usepackage{enumitem}
\usepackage[hidelinks]{hyperref}

\title{Stationary weighted Schiffer: perturbation on the ball}
\author{}
\date{\today}

\newcommand{\R}{\mathbb{R}}
\newcommand{\Om}{\Omega}
\newcommand{\eps}{\varepsilon}
\newcommand{\1}{\mathbf 1}
\newcommand{\nor}{\partial_\nu}
\newcommand{\tr}{\mathrm{tr}}

\newtheorem{lem}{Lemma}
\newtheorem{prop}[lem]{Proposition}
\newtheorem{cor}[lem]{Corollary}
\newtheorem{thm}[lem]{Theorem}
\newtheorem{rem}[lem]{Remark}
\newtheorem{conj}[lem]{Conjecture}
\newtheorem{defn}[lem]{Definition}

\begin{document}
\maketitle

\section{Where we are}

The note \texttt{direction2-parabolic-to-classical.tex} proved that
the lateral parabolic Schiffer condition for a weighted pair
$(\Om,V)$ reduces to classical Schiffer for the unweighted Laplacian
on the base. That closes the parabolic angle on Direction 2 (and
makes precise what the de Dios Pont barrel limit can and cannot do).
The residual question is the \emph{stationary} weighted Schiffer
problem:

\begin{defn}[Stationary weighted Schiffer pair]\label{def:wsp}
A smooth convex pair $(\Om,V)$ with eigenpair $(\phi,\lambda)$
is a \emph{stationary weighted Schiffer pair} if
\begin{equation}\label{eq:swsp}
    L_{\Om,V}\phi := -\Delta\phi+\nabla V\cdot\nabla\phi=\lambda\phi
    \ \text{in }\Om,\quad
    \nor\phi=0\ \text{on }\partial\Om,\quad
    \phi|_{\partial\Om}\equiv c,\quad \phi\not\equiv 0.
\end{equation}
\end{defn}

The ball with any radial convex $V$ produces a countable family of
such pairs (one per radial Langevin eigenvalue). The natural
extension of the Schiffer conjecture is:

\begin{conj}[Log-concave Schiffer]\label{conj:lcs}
If $(\Om,V)$ is a smooth convex stationary weighted Schiffer pair,
then $\Om$ is a ball and $V$ is radial (with respect to the centre of
$\Om$).
\end{conj}

This generalizes the classical Schiffer conjecture (recovered for
$V\equiv 0$) and is the right target for Direction 2 with the
log-concave compactification.

Below we attack the conjecture perturbatively at the unweighted
endpoint $V=0$ on the ball, and identify the precise codimension of
the Schiffer-preserving directions of $V$. The output is:

\begin{itemize}[leftmargin=*]
    \item Schiffer-preserving first-order deformations $V_1$ on the
    ball form an infinite-codimensional but still infinite-dimensional
    linear subspace; the codimension is enumerated by an explicit
    family of linear functionals $\mathcal{L}_{\ell,m}$ indexed by
    nonzero spherical harmonics $(\ell,m)$, each of which is a
    one-dimensional Green-function moment of $V_{1,\ell,m}$ against
    the radial profile $\phi_0'$.
    \item The radial directions $V_1=V_1(r)$ are all in this kernel,
    matching the known continuum of radial weighted Schiffer pairs.
    \item The infinitesimal log-concave Schiffer rigidity question on
    the ball reduces to: \emph{does every first-order Schiffer-preserving
    deformation extend to all orders, modulo $O(n)$ symmetries?} A
    positive answer would say the perturbation is forced into the
    radial direction.
\end{itemize}

\section{Setup}

Throughout this note $\Om=B=B_R\subset\R^n$, $n\ge 2$. Fix a radial
Neumann eigenfunction
\[
    \phi_0(r) \quad\text{with}\quad
    -\Delta\phi_0=\lambda_0\phi_0\ \text{in }B,\quad
    \phi_0'(R)=0,
\]
i.e., $\phi_0$ is a Bessel-type radial eigenfunction at a radial
Neumann eigenvalue $\lambda_0$. By radiality $\phi_0(R)=c_0$ is
automatically a constant boundary value, so $(B,0,\phi_0,\lambda_0)$
is a (trivial) weighted Schiffer pair.

Perturb with $V_\eps=\eps V_1$, $V_1\in C^\infty(\overline B)$, and
expand $\phi_\eps=\phi_0+\eps\phi_1+O(\eps^2)$,
$\lambda_\eps=\lambda_0+\eps\lambda_1+O(\eps^2)$.

\section{First-order eigenvalue equation}

\begin{lem}\label{lem:order1}
$(\phi_1,\lambda_1)$ satisfies
\begin{equation}\label{eq:firstorder}
    (-\Delta-\lambda_0)\phi_1=\lambda_1\phi_0-\nabla V_1\cdot\nabla\phi_0
    \ \text{in }B,\qquad
    \nor\phi_1=0\ \text{on }\partial B,
\end{equation}
and the Fredholm condition gives
\begin{equation}\label{eq:lambda1}
    \lambda_1=\frac{1}{\|\phi_0\|^2_{L^2(B)}}
    \int_B(\nabla V_1\cdot\nabla\phi_0)\,\phi_0\,dx
    =\frac{1}{2\|\phi_0\|^2_{L^2(B)}}
    \int_B \nabla V_1\cdot\nabla(\phi_0^2)\,dx.
\end{equation}
\end{lem}

\begin{proof}
Substitute the expansion into $L_{B,\eps V_1}\phi_\eps=\lambda_\eps\phi_\eps$
and collect the $O(\eps)$ term. The Neumann BC at first order is
$\nor\phi_1=0$ since both $\phi_0$ and $\phi_\eps$ have $\nor=0$ on
$\partial B$. Self-adjointness of $-\Delta-\lambda_0$ with Neumann data
on its kernel $\mathrm{span}\{\phi_0\}$ yields \eqref{eq:lambda1}.
\end{proof}

\section{Decomposition into spherical harmonics}

Let $\{Y_{\ell,m}\}_{\ell\ge 0,m}$ be the standard orthonormal basis
of $L^2(S^{n-1})$ with $-\Delta_{S^{n-1}}Y_{\ell,m}=\ell(\ell+n-2)Y_{\ell,m}$.
For $\ell=0$ take $Y_{0,1}\equiv |S^{n-1}|^{-1/2}$.
Decompose
\[
    V_1(x)=\sum_{\ell,m}V_{1,\ell,m}(r)\,Y_{\ell,m}(\hat x),\qquad
    \phi_1(x)=\sum_{\ell,m}\Phi_{\ell,m}(r)\,Y_{\ell,m}(\hat x).
\]
Since $\phi_0$ is radial, $\nabla V_1\cdot\nabla\phi_0=\phi_0'(r)\partial_r V_1
=\phi_0'(r)\sum_{\ell,m}V_{1,\ell,m}'(r)Y_{\ell,m}(\hat x)$.
Write the radial Laplacian in the $\ell$-th sector as
\[
    \Delta_r^{(\ell)}u(r):=u''(r)+\tfrac{n-1}{r}u'(r)
        -\tfrac{\ell(\ell+n-2)}{r^2}u(r).
\]
Equation \eqref{eq:firstorder} separates:
\begin{equation}\label{eq:ellsector}
    (-\Delta_r^{(\ell)}-\lambda_0)\Phi_{\ell,m}(r)
    =\delta_{\ell,0}\,\lambda_1\,\phi_0(r)\,\delta_{m,1}|S^{n-1}|^{1/2}
        -\phi_0'(r)\,V_{1,\ell,m}'(r),
\end{equation}
with $\Phi_{\ell,m}'(R)=0$ and regularity at $r=0$.

\section{The Schiffer-preserving constraint at first order}

The first-order weighted Schiffer condition is
$\phi_\eps|_{\partial B}\equiv c_\eps$, which gives $\phi_1(R\hat x)=c_1$
for every unit vector $\hat x\in S^{n-1}$. In spherical-harmonic
modes:
\begin{equation}\label{eq:schifferconstraint}
    \Phi_{0,1}(R)=c_1\,|S^{n-1}|^{1/2},\qquad
    \Phi_{\ell,m}(R)=0\quad\text{for all }(\ell,m)\text{ with }\ell\ge 1.
\end{equation}
The $\ell=0$ equation just adjusts the boundary constant by $\eps c_1$,
which is allowed (the Schiffer condition only demands constancy in
the angular variable, not a prescribed value); so $\ell=0$ imposes no
constraint on $V_1$. The genuine constraint is on the $\ell\ge 1$
modes.

\begin{prop}[First-order rigidity functional]\label{prop:functional}
For each $\ell\ge 1$ and each $m$, the constraint
$\Phi_{\ell,m}(R)=0$ is equivalent to the vanishing of the linear
functional
\begin{equation}\label{eq:functional}
    \mathcal{L}_{\ell,m}(V_1)
    :=\int_0^R G_\ell(R,r')\,\phi_0'(r')\,V_{1,\ell,m}'(r')\,dr',
\end{equation}
where $G_\ell(R,r')$ is the Green's function on $[0,R]$ for the
operator $-\Delta_r^{(\ell)}-\lambda_0$ with regular boundary
behaviour at $r=0$ and Neumann condition $u'(R)=0$, evaluated at
$r=R$.
\end{prop}

\begin{proof}
For $\ell\ge 1$ the right-hand side of \eqref{eq:ellsector} is just
$-\phi_0'(r)V_{1,\ell,m}'(r)$. Solve by convolution against the
Green's function and read off the value at $r=R$.
\end{proof}

\begin{rem}[Existence and non-degeneracy of $G_\ell$]\label{rem:gell}
For generic radial eigenvalues $\lambda_0$, $\lambda_0$ is not a
Neumann eigenvalue of $-\Delta_r^{(\ell)}$ for $\ell\ge 1$, so
$G_\ell$ exists. When $\lambda_0$ accidentally lies in the
$\ell$-sector Neumann spectrum a finite-codimensional Fredholm
adjustment is required, but generically does not occur for the
radial Bessel eigenvalues used here.
\end{rem}

\begin{cor}[The first-order kernel]\label{cor:kernel}
The space of $V_1\in C^\infty(\overline B)$ such that
$(B,\eps V_1)$ admits a first-order Schiffer-preserving extension of
$(\phi_0,\lambda_0)$ is the linear subspace
\begin{equation}\label{eq:kernel}
    \mathcal{K}:=\bigcap_{\ell\ge 1,\,m}\ker\mathcal{L}_{\ell,m}\subset C^\infty(\overline B).
\end{equation}
\end{cor}

\section{Radial perturbations and the size of $\mathcal{K}$}

\begin{cor}[Radial directions are admissible]\label{cor:radial}
Every radial $V_1=V_1(r)$ lies in $\mathcal{K}$.
\end{cor}

\begin{proof}
For radial $V_1$, $V_{1,\ell,m}\equiv 0$ for all $(\ell,m)$ with
$\ell\ge 1$, so $\mathcal{L}_{\ell,m}(V_1)=0$ trivially.
\end{proof}

\begin{cor}[Beyond radial]\label{cor:beyond}
$\mathcal{K}$ strictly contains the radial subspace. In particular,
for every $(\ell,m)$ with $\ell\ge 1$, the kernel
$\ker\mathcal{L}_{\ell,m}$ inside the angular sector $V_{1,\ell,m}$
is infinite-codimension-one in $C^\infty([0,R])$; the joint
intersection over $(\ell,m)$ remains infinite-dimensional.
\end{cor}

\begin{proof}
$\mathcal{L}_{\ell,m}$ is a single nonzero linear functional on the
infinite-dimensional space of radial profiles $V_{1,\ell,m}$, so its
kernel is of codimension 1 in that sector. Different $(\ell,m)$
sectors are independent, so the joint kernel intersects each sector
in a codimension-1 subspace.
\end{proof}

This says the first-order Schiffer-preserving perturbations are an
infinite-dimensional family that strictly contains the radial
directions. If Conjecture~\ref{conj:lcs} is true, all but the radial
directions must be ruled out by higher-order obstructions or by
removing $O(n)$-trivial directions.

\section{$O(n)$-trivial directions}

The $O(n)$-action on $B$ produces a finite-dimensional family of
``trivial'' perturbations that do not change the pair up to symmetry.
Each angular harmonic $Y_{\ell,m}$ for $\ell=1$ is generated by an
infinitesimal translation of the centre of $B$ (a rigid motion):
specifically, $V_1(x)=x_j$ corresponds to translating the centre of
the sphere. Including these in the analysis: the $\ell=1$ sector of
$\mathcal{K}$ is exactly the span of $\{x_j\cdot g(r):g\in
\ker\mathcal L_{1,j}\}$, and modulo translations of the centre this
sector is reduced by $n$ dimensions.

Higher $\ell\ge 2$ harmonics are not generated by rigid motions, so
their kernel $\ker\mathcal{L}_{\ell,m}$ is genuinely
infinite-dimensional beyond symmetries.

This is the key observation: \emph{infinitesimal log-concave Schiffer
rigidity on the ball is equivalent to: for every $\ell\ge 2$ and every
$V_{1,\ell,m}\in\ker\mathcal{L}_{\ell,m}\setminus\{0\}$, the
formal Lyapunov--Schmidt series for $(\phi_\eps,\lambda_\eps,V_\eps)$
fails to converge to a genuine weighted Schiffer pair at some finite
order $\ge 2$.}

\section{The second-order test}

The natural follow-up computation is the second-order Schiffer
condition. Write $\phi_\eps=\phi_0+\eps\phi_1+\eps^2\phi_2+\cdots$.
The $O(\eps^2)$ equation is
\begin{equation}\label{eq:order2eqn}
    (-\Delta-\lambda_0)\phi_2
    =\lambda_2\phi_0+\lambda_1\phi_1-\nabla V_1\cdot\nabla\phi_1.
\end{equation}
At order 2, the Schiffer condition $\phi_2(R\hat x)=c_2$ gives, for
each $(\ell,m)$ with $\ell\ge 1$,
\begin{equation}\label{eq:order2constraint}
    \int_0^R G_\ell(R,r')\bigl[\lambda_1\Phi_{\ell,m}(r')
        -(\nabla V_1\cdot\nabla\phi_1)_{\ell,m}(r')\bigr]dr' = 0,
\end{equation}
which involves the angular-product structure
$(V_1\cdot\phi_1)_{\ell,m}$. This is a quadratic functional in
$V_1$. The key analytic question is:

\begin{quote}
\emph{Is there a non-radial $V_1\in\mathcal{K}$ such that
\eqref{eq:order2constraint} holds for all $(\ell,m)$ with $\ell\ge
1$, and similarly all higher-order constraints?}
\end{quote}

A positive answer produces a genuine non-radial branch of weighted
Schiffer pairs on the ball, falsifying Conjecture~\ref{conj:lcs} at
the perturbative level. A negative answer (which the conjecture
predicts) requires showing that the quadratic constraint
\eqref{eq:order2constraint} together with the linear
$\mathcal{L}_{\ell,m}=0$ system has only the radial solutions.

\section{Next steps}

\begin{enumerate}[leftmargin=*]
    \item \textbf{Compute $G_\ell(R,r')$ explicitly} for the disc
    ($n=2$) and ball ($n=3$) using Bessel and trigonometric
    representations, and check whether $\mathcal{L}_{\ell,m}$ has
    sign-definite Green-function kernel against $\phi_0'$. If
    so, $\mathcal{L}_{\ell,m}(V_{1,\ell,m})$ controls
    $V_{1,\ell,m}'$ moments in a useful way.
    \item \textbf{Check $\ell=1$ explicitly.} The $\ell=1$ sector is
    where rigid-motion symmetries live; the corresponding kernel
    should be exactly the translation directions modulo symmetry. A
    direct check confirms whether the first-order analysis ``sees''
    the rigid-motion freedom correctly.
    \item \textbf{Quadratic test for $\ell=2$.} The smallest
    non-trivial angular harmonic at which the conjecture would be
    tested is $\ell=2$. Carry out
    \eqref{eq:order2constraint} for $V_1=V_{1,2,m}(r)Y_{2,m}$ in
    $\ker\mathcal{L}_{2,m}$ and determine whether any non-zero such
    $V_1$ is consistent.
    \item \textbf{Connect to known classical results.} The
    second-order test is, when $V_1$ is taken instead as an
    infinitesimal boundary perturbation of the disc (not a potential
    perturbation), the standard Lyapunov--Schmidt analysis for
    classical Schiffer near the disc (Direction 3 of the original
    notes). The relationship between these two perturbations should
    illuminate which features of the Schiffer obstruction come from
    geometry of $\partial\Om$ versus from the potential $V$.
\end{enumerate}

\end{document}
